% This section is intended for inclusion in the consolidated English appendix.
% Bibliography keys used below: vdER, Shastri, Craighero, Masuda.
% All conclusions depend on the complete rectifier, not merely on a parameter lift.
\section{Which named conjectures are affected by rectification}

In this section, ``refutation'' means a complete mathematical argument
with exact certificates, presented for verification. External peer review
and priority in the literature are not claimed.

Set $K=\mathbb Q(\sqrt{-3})$ and
\[
\gamma_T(u)=(f(u),g(u),h(u))
=(u^3-3u,u^4-4u^2,u^5-10u).
\]
Three equivalent conjectures about a specific object associated with this
embedding are formulated in \cite[pp.~1319--1320]{vdER}.
They assert, respectively, that the kernel of a certain LND, the quotient
by its slice, and the extended Rees algebra of the curve ideal are not
polynomial rings. The last formulation is called Asanuma's Conjecture there.
The argument below, applied to the certified rectifying map over $K$,
negates all three formulations over $\mathbb C$.

\begin{proposition}[Explicit coordinates for the associated LND]
Let $k$ be a field of characteristic zero,
$F=(P,Q,R_3)\in\operatorname{Aut}_k(\mathbb A^3)$,
$G=F^{-1}$, $F\circ\gamma=(u,0,0)$ and $0\ne b(t)\in k[t]$.
For
\[
A=k[t,u,X_1,X_2,X_3],\qquad
D_b=b(t)\partial_u+\sum_i f_i'(u)\partial_{X_i}
\]
there is a polynomial coordinate system
$A=k[t,p,q,r,s]$ in which $D_b=\partial_s$.
\end{proposition}

\begin{proof}
Set $V_i=f_i(u)-b(t)X_i$ and
\[
p=P(V),\qquad q=\frac{Q(V)}{b(t)},\qquad
r=\frac{R_3(V)}{b(t)},\qquad
s=\frac{u-P(V)}{b(t)}.
\]
Every numerator is exactly divisible by $b(t)$, since
$V\equiv\gamma(u)\pmod{b(t)}$. The inverse substitution is given by
\[
u=p+b(t)s,\qquad
X_i=\frac{f_i(p+b(t)s)-G_i(p,b(t)q,b(t)r)}{b(t)}.
\]
It is also polynomial because $G(p,0,0)=\gamma(p)$.
Both compositions are the identity after inverting $b(t)$; injectivity
of localization proves the same identities in the original ring. Finally,
$D_b(V_i)=0$, so $D_b(p)=D_b(q)=D_b(r)=0$ and $D_b(s)=1$.
\end{proof}

For the Shastri model with $b(t)=t$, this gives
\[
\ker D=K[t,p,q,r]\simeq K^{[4]}.
\]
After extension of scalars, this contradicts the nonpolynomiality
asserted in Conjecture~5.1 \cite{vdER}.

Shastri's original parameter polynomial is
\[
F_0(x,y,z)=yz-x^3-5xy+2z-7x,
\]
and the slice in Conjecture~5.2 is
$\sigma=(u-F_0(V))/t$. Since $D(\sigma)=D(s)=1$,
the difference $\sigma-s$ belongs to $B=\ker D$.
Consequently, $A=B[\sigma]$; in particular,
\[
A/(\sigma-c)\simeq B\simeq K^{[4]}\qquad(c\in K).
\]
This negates Conjecture~5.2 after passing to $\mathbb C$.
Nevertheless, $F_0$ itself is not a coordinate of $\mathbb A^3$:
its gradient vanishes at $(1,-2,5)$ and $(-1,-2,-5)$.
The fact that $F_0$ is not a coordinate is compatible with $\sigma$ being
a coordinate, since these are different polynomials in different rings.

Finally, with $R_0=K[x,y,z]$ and $I=I_{\gamma_T}$, we have
$R_0=K[P,Q,R_3]$, $I=(Q,R_3)$, and hence
\[
R_0[t,I/t]=K[t,P,Q/t,R_3/t]\simeq K^{[4]}.
\]
Here $R_0[t,I/t]\subset R_0[t,t^{-1}]$ is precisely the extended
Rees algebra occurring in Conjecture~5.3. This negates the third
formulation; it is not a result about arbitrary Rees algebras.

Thus one particular candidate is eliminated in three related presentations.
This does not solve the general cancellation problem. Neither does it solve
the general rectification problem for $\mathbb A^1\hookrightarrow
\mathbb A^3$ or the general conjecture that nonrectifiable space lines exist.

\paragraph{Historical distinctions.}
The conjectured complex nonrectifiability of this specific $\gamma_T$
is attributed to Shastri in \cite[p.~1319]{vdER}; the supplied rectifying
map refutes it. However, Questions Q.1 and Q.2 in
\cite[p.~12]{Shastri} were already solved by Shastri: they concern
polynomial representations of real knots whose complexifications are
closed embeddings, rather than their complex rectifiability.
Special degree-based conjectures of Abhyankar had earlier refutations,
as described in \cite[Introduction, pp.~115--116]{Craighero}.

\paragraph{Fields and actions.}
The resulting coordinates are defined over $K$ and hence over $\mathbb C$;
descent to $\mathbb Q$ or $\mathbb R$ has not been proved.
When $b=t$, the natural action
$t\mapsto\lambda t$, $u\mapsto u$, $X_i\mapsto\lambda^{-1}X_i$
is linear in the kernel coordinates, with weights $(1,0,-1,-1)$.
This is a specific linearization. It does not imply statements about all
group actions. In the reverse implications toward the embedding, the
grading is essential, as in \cite{Masuda}; it cannot be replaced merely
by abstract polynomiality of the kernel or by $\operatorname{ML}(B)=k$.

\subsection{Full ML, relative ML and the filtration}
For a rectified model $B=k[t,p,q,r]$, the four coordinate
derivations give
\[
\ML(B)=k,\qquad \ML^*(B)=k,\qquad \operatorname{Dk}(B)=B.
\]
Here $\ML^*$ uses only LNDs with slices, and $\operatorname{Dk}$
is the algebra generated by the kernels of all nonzero LNDs.
For $D_b=\partial_s$, we have $\pl(D_b)=B$,
$\ker D_b^{n+1}=\bigoplus_{j=0}^n Bs^j$ and
$D_b^n(\ker D_b^{n+1})=B$.
This is consistent with a large minimum slice degree in the original
coordinates: these are different characteristics.

For an arbitrary embedding whose rectifying map is unknown, the following
dichotomy still holds:
\[
\ML(B_b)\in\{k,k[t]\},\qquad B_b=\ker D_b.
\]
Indeed, the derivations $\partial_{X_i}$ commute with $D_b$ and restrict
to $B_b$. The intersection of their kernels there is
$B_b\cap k[t,u]=k[t]$. If a nonconstant $a(t)$ belongs to the full ML,
then for every LND $\Delta$ on $B_b$, the equality
$0=\Delta a(t)=a'(t)\Delta t$ in a characteristic-zero domain
implies $\Delta t=0$. Thus all of $k[t]$ belongs to ML.

Therefore, $\ML(B_b)=k[t]$ would obstruct rectification.
However, $\ML(B_b)=k$ proves neither polynomiality of $B_b$
nor rectifiability. If only LNDs fixing $t$ are allowed from the outset,
the relative ML equals $k[t]$ even in the polynomial case.
For the new kernels associated with $H(3,7,8)$ and $H(4,5,7)$,
the relevant branch of the dichotomy has not yet been determined.

\subsection{What must be retained for the reverse implication}
Let $R=k[x,y,z]$, $I=I_\gamma$,
$B_\mu=R[t,I/t^\mu]$, $\mu\ge1$, $\wt(t)=1$, $\wt(R)=0$.
The embedding is rectifiable if and only if there are homogeneous
coordinates
\[
B_\mu=k[t,p,q,r],\qquad
\wt(t,p,q,r)=(1,0,-\mu,-\mu)
\]
with \emph{the same distinguished element $t$}. The reverse implication
follows from
\[
\begin{gathered}
(B_\mu)_0=R=k[p,t^\mu q,t^\mu r],\\
I=t^\mu(B_\mu)_{-\mu}=R\cap tB_\mu=(t^\mu q,t^\mu r).
\end{gathered}
\]
The restriction of $p$ to the curve is an affine parameter and can
be normalized to $u$.

An equivalent formulation is the existence of a homogeneous LND $E$
of degree $-1$ on $B_\mu$ with $E(t)=1$. The forward implication
is supplied by $\partial_t$. Conversely, the slice theorem and the finite projection
\[
\rho(b)=\sum_{j\ge0}\frac{(-t)^j}{j!}E^j(b)
\]
give $B_\mu=(\ker E)[t]$ and a graded isomorphism
$\ker E\simeq B_\mu/(t)$. This special fibre is polynomial with weights
$(0,-\mu,-\mu)$ by the normal-basis construction in Section~\ref{sec:formal}.
Thus the required homogeneous coordinates lift.

For an LND $E$, without assuming homogeneity, the condition $E(t)=1$
still gives abstract polynomiality of $B_\mu$, but the preceding proof
no longer recovers the embedding.
For an LND $E$ satisfying only $E(t)\ne0$, one obtains only $\ML(B_\mu)=k$.
One cannot take an intermediate homogeneous component of an arbitrary LND
and declare it locally nilpotent without proof.

\paragraph{Direct connection with the planar theorem.}
If the ideal $I$ contains a polynomial coordinate, the curve can be placed
in a coordinate plane and the Abhyankar--Moh--Suzuki theorem applied.
If a parameter polynomial $P\circ\gamma=u$ is itself a coordinate, then
with $R=k[P,U,V]$ the required rectifying map is immediately
$(P,U-U(\gamma(P)),V-V(\gamma(P)))$.
A parameter polynomial that is not necessarily a coordinate is
insufficient for either argument.
